\chapter{An\'alisis temporal y costes de desarrollo}\label{anatemporal}
% tfg-floatbarrier-auto
\makeatletter
\@ifundefined{tfgfloatbarrier}{
  \def\tfgfloatbarrier{1}
  \let\tfgorigfigure\figure
  \let\tfgendorigfigure\endfigure
  \renewenvironment{figure}{\tfgorigfigure}{\tfgendorigfigure\FloatBarrier}
  \let\tfgorigtable\table
  \let\tfgendorigtable\endtable
  \renewenvironment{table}{\tfgorigtable}{\tfgendorigtable\FloatBarrier}
}{}
\makeatother

\FloatBarrier
\section{An\'alisis temporal}
\subsection{Metodolog\'ia aplicada y roles del proyecto}

La ejecuci\'on del TFG se ha organizado con un enfoque \textbf{\'agil adaptado a contexto acad\'emico},
tomando Scrum como marco de referencia ligero.
No se replica un proceso empresarial completo,
pero s\'i se mantienen los elementos que aportan mayor valor de control:
iteraciones cortas, objetivos por sprint, revisi\'on frecuente y replanificaci\'on basada en evidencia.

Los roles de trabajo utilizados han sido el \textbf{equipo de desarrollo}, formado por dos autores con responsabilidad compartida en backend, frontend, pruebas y documentaci\'on t\'ecnica; el \textbf{tutor del TFG}, con funci\'on de supervisi\'on funcional y t\'ecnica para validar hitos, alcance y calidad de entregables; y los \textbf{stakeholders acad\'emicos}, como referencia de necesidades operativas reales en docencia, entrega y evaluaci\'on. Esta combinaci\'on permite equilibrar rigor t\'ecnico con alineaci\'on a uso docente.

La duraci\'on de iteraci\'on se fij\'o en dos semanas,
con entregables funcionales incrementales y una revisi\'on de cierre por sprint.

\subsection{Herramientas de gesti\'on y seguimiento}

La gesti\'on del trabajo se ha apoyado en artefactos versionados del propio repositorio y en herramientas de planificaci\'on integradas con el flujo de desarrollo. En concreto, se utilizaron \textbf{Product Backlog} y \textbf{Sprint Backlog} para priorizaci\'on y control de alcance, \textbf{GitHub Issues/Projects} para trazabilidad de tareas y decisiones de ejecuci\'on, y un \textbf{registro horario por sprint} que contrasta capacidad prevista con dedicaci\'on real. Esta triangulaci\'on permite justificar desviaciones de forma auditable.

Este esquema permite justificar desviaciones de forma auditable,
conectando horas, funcionalidades y evidencia de validaci\'on.

\subsection{Marco de planificaci\'on}

La planificaci\'on temporal del proyecto se ha definido siguiendo un enfoque incremental
basado en \emph{sprints} quincenales, con revisiones peri\'odicas del alcance.
El objetivo de este enfoque no es \'unicamente repartir tareas, sino reducir riesgo t\'ecnico
de forma temprana y asegurar que cada incremento entregue valor funcional verificable.

El marco de trabajo adoptado combina planificaci\'on por objetivos (cada sprint se orienta a un bloque funcional medible), control continuo con seguimiento semanal de avance, calidad y bloqueos, y replanificaci\'on expl\'icita cuando cambian supuestos. Esta mezcla aporta disciplina sin perder capacidad de ajuste.

\subsection{Hip\'otesis de capacidad y calendario base}

Los supuestos de capacidad empleados para el plan temporal fueron un equipo de desarrollo de \textbf{2 personas}, con dedicaci\'on estimada de \textbf{30 horas por persona y semana}, lo que fija una capacidad total semanal de \textbf{60 horas}. Se definieron sprints de \textbf{2 semanas}, con una capacidad por sprint de \textbf{120 horas}. Estos supuestos hacen expl\'icita la base del calendario y permiten medir desviaciones.

Sobre esa base, el proyecto se estructur\'o en \textbf{6 iteraciones} (Sprint 0 + Sprint 1 a Sprint 5),
con un esfuerzo total planificado de \textbf{660 horas}.
Se a\~nadi\'o un Sprint 5 de cierre, con menor carga (60 h) y calendario extendido
hasta el 12 de mayo para revisi\'on final, correcci\'on de incidencias y ajuste documental.

\subsection{Planificaci\'on macro por sprints}

La Tabla~\ref{tab:plan-macro-sprints} resume la planificaci\'on temporal de alto nivel.

\begin{table}[!htbp]
\centering
\begin{tabular}{|P{1.3cm}|P{2.8cm}|P{1.8cm}|P{5.8cm}|}
\hline
\textbf{Sprint} & \textbf{Periodo} & \textbf{Horas} & \textbf{Objetivo principal} \\
\hline
S0 & Sep 2025 - 22 Feb 2026 & 120 & Documentaci\'on y an\'alisis inicial: ERS, DAS, casos de uso, modelo de datos y backlog. \\
\hline
S1 & 17 Feb - 2 Mar 2026 & 120 & Infraestructura base, arquitectura backend/frontend, autenticaci\'on, m\'odulos n\'ucleo y despliegue inicial. \\
\hline
S2 & 3 Mar - 16 Mar 2026 & 120 & Funcionalidad de entregables y entregas, mejora del flujo de estudiante/profesor e integraciones previstas. \\
\hline
S3 & 17 Mar - 30 Mar 2026 & 120 & Evaluaci\'on y feedback, mejoras de usabilidad para profesorado, refuerzo de pruebas e inicio formal de memoria. \\
\hline
S4 & 31 Mar - 13 Abr 2026 & 120 & Cierre t\'ecnico, pruebas finales, documentaci\'on completa, pulido del producto y preparaci\'on de defensa. \\
\hline
S5 & 14 Abr - 12 May 2026 & 60 & Revisi\'on final, correcci\'on de incidencias, ajustes de memoria y cierre documental. \\
\hline
\end{tabular}
\caption{Planificaci\'on temporal macro del proyecto}
\label{tab:plan-macro-sprints}
\end{table}

\subsection{Planificaci\'on micro y distribuci\'on del esfuerzo}

Adem\'as de la visi\'on macro, cada sprint se descompuso en l\'ineas de trabajo concretas
con asignaci\'on horaria por responsable. El criterio aplicado fue identificar primero las tareas cr\'iticas de MVP, reservar capacidad para pruebas en cada iteraci\'on y mantener un \emph{buffer} expl\'icito para incidencias. Esta pauta reduce el riesgo de concentrar la validaci\'on en la fase final.

En t\'erminos acumulados, la distribuci\'on de esfuerzo se ha orientado a construir la base arquitect\'onica y de seguridad en fases tempranas, desarrollar funcionalidad de negocio en fases intermedias y concentrar validaci\'on de calidad, documentaci\'on y cierre en la fase final. Esta secuencia protege la estabilidad del sistema a medida que crece el alcance.

Esta distribuci\'on evita concentrar el testing al final y reduce el riesgo de deuda t\'ecnica
de cierre, un problema habitual en proyectos con calendario acad\'emico estricto.

\subsection{Hitos de control del proyecto}

Para asegurar gobernanza temporal se definieron hitos de control: \textbf{H1} (cierre de requisitos y arquitectura con ERS/DAS validados), \textbf{H2} (disponibilidad de la base t\'ecnica del producto en backend, frontend y seguridad), \textbf{H3} (primer flujo operativo extremo a extremo actividad-entrega), \textbf{H4} (despliegue operativo y pipelines CI/CD estables), \textbf{H5} (incorporaci\'on del m\'odulo de evaluaci\'on y feedback), \textbf{H6} (estabilizaci\'on de calidad con pruebas autom\'aticas y mutaci\'on), \textbf{H7} (cierre de memoria t\'ecnica y validaci\'on documental) y \textbf{H8} (preparaci\'on de defensa y cierre final del proyecto). Esta secuencia introduce puntos de control verificables en momentos cr\'iticos.

\subsection{Seguimiento del avance y m\'etricas de control}

El seguimiento temporal se realiz\'o con m\'etricas sencillas: horas planificadas frente a horas completadas por sprint, porcentaje acumulado de ejecuci\'on respecto a las 660 horas totales, velocidad de entrega por iteraci\'on (PBIs o tareas completadas) e indicadores de calidad como condici\'on de cierre (tests en verde, cobertura y estabilidad de pipelines). Estas m\'etricas conectan avance, alcance y calidad en una misma lectura.

Como referencia de control acumulado, tras S0 se ejecutaron 120 h (18\%), tras S1 se alcanzaron 240 h (36\%), el objetivo tras S2 fue 360 h (55\%), tras S3 480 h (73\%), tras S4 600 h (91\%) y tras S5 660 h (100\%). Esta trayectoria permite comprobar si la ejecuci\'on se mantiene alineada con el plan.

\subsection{Control de desviaciones con valor ganado simplificado}

Para evaluar el progreso del proyecto de forma objetiva,
se ha utilizado una adaptaci\'on del m\'etodo de an\'alisis de valor ganado (EVM)
acorde al alcance de un trabajo acad\'emico.

Se consideran tres magnitudes: \textbf{PV} (\emph{Planned Value}), como horas planificadas acumuladas; \textbf{EV} (\emph{Earned Value}), como horas equivalentes al trabajo realmente completado; y \textbf{AC} (\emph{Actual Cost}), como horas reales consumidas. Esta base permite interpretar desviaciones de manera objetiva.

Los indicadores de interpretaci\'on utilizados son:

\[
SPI = \frac{EV}{PV}
\qquad
CPI = \frac{EV}{AC}
\]

donde un valor cercano a 1 indica equilibrio entre plan y ejecuci\'on.

\begin{table}[!htbp]
\centering
\begin{tabular}{|P{2cm}|P{2.2cm}|P{2.2cm}|P{2.2cm}|P{1.5cm}|P{1.5cm}|}
\hline
\textbf{Corte} & \textbf{PV (h)} & \textbf{EV (h)} & \textbf{AC (h)} & \textbf{SPI} & \textbf{CPI} \\
\hline
Fin S1 & 240 & 236 & 242 & 0.98 & 0.98 \\
\hline
Fin S2 & 360 & 354 & 366 & 0.98 & 0.97 \\
\hline
Fin S3 & 480 & 474 & 486 & 0.99 & 0.98 \\
\hline
Fin S4 (objetivo) & 600 & 600 & 600 & 1.00 & 1.00 \\
\hline
Fin S5 (objetivo) & 660 & 660 & 660 & 1.00 & 1.00 \\
\hline
\end{tabular}
\caption{Seguimiento temporal con valor ganado simplificado}
\label{tab:valor-ganado}
\end{table}

La lectura de estos valores indica desviaciones controladas,
sin deriva grave de calendario,
lo que valida la eficacia de la replanificaci\'on incremental.

\subsection{Desviaciones y replanificaci\'on}

El plan no permaneci\'o est\'atico durante todo el proyecto. Se registraron revisiones
de planificaci\'on con trazabilidad de versiones, destacando el ajuste del modelo de capacidad (de una planificaci\'on inicial m\'as fragmentada a una estructura de sprints de 120 h), la recalibraci\'on de fechas de S1-S4 para alinear ejecuci\'on t\'ecnica y calendario acad\'emico, y la incorporaci\'on de tareas documentales en S3 para evitar cuello de botella en S4. Estas revisiones preservan el objetivo global sin perder consistencia temporal.

La replanificaci\'on se adopt\'o con criterio de continuidad: no alter\'o el objetivo global de 660 h,
pero mejor\'o la probabilidad de cumplimiento de hitos y redujo el riesgo de cierre tard\'io.

\subsection{Lecciones aprendidas de planificaci\'on}

Del seguimiento temporal se extraen tres lecciones principales. Iniciar testing temprano reduce retrabajo porque distribuir pruebas por sprint evita acumulaci\'on de defectos al final. Documentar en paralelo reduce riesgo de cierre, ya que incorporar memoria desde S3 disminuye el cuello de botella de la fase final. Planificar un \emph{buffer} expl\'icito mejora estabilidad al absorber incidencias sin romper hitos cr\'iticos.

\subsection{Riesgos temporales y mitigaci\'on}

Durante la planificaci\'on se identificaron riesgos de calendario relevantes.
La Tabla~\ref{tab:riesgos-temporales} resume los principales.

\begin{table}[!htbp]
\centering
\begin{tabular}{|P{4.9cm}|P{2cm}|P{2cm}|P{4.0cm}|}
\hline
\textbf{Riesgo} & \textbf{Prob.} & \textbf{Impacto} & \textbf{Mitigaci\'on} \\
\hline
Complejidad UX en panel docente de evaluaci\'on & Media & Alto & Prototipado temprano y validaci\'on incremental por sprint. \\
\hline
Sobrecarga de redacci\'on de memoria en fase final & Alta & Alto & Inicio de memoria en S3 y redacci\'on paralela al desarrollo. \\
\hline
Incidencias de integraci\'on frontend-backend & Media & Medio & Pruebas de integraci\'on continuas y criterios de cierre por calidad. \\
\hline
Variabilidad de disponibilidad semanal del equipo & Media & Alto & Reserva de buffer de contingencia por iteraci\'on. \\
\hline
\end{tabular}
\caption{Riesgos temporales y acciones de mitigaci\'on}
\label{tab:riesgos-temporales}
\end{table}

\subsection{Conclusi\'on temporal}

Desde el punto de vista temporal, la estrategia adoptada combina rigor de planificaci\'on
y adaptabilidad operativa. La descomposici\'on por sprints, el seguimiento con m\'etricas de avance
y la replanificaci\'on controlada permiten sostener una ejecuci\'on realista,
alineada con los requisitos del TFG y con las restricciones del calendario acad\'emico.

\FloatBarrier
\section{Costes de desarrollo}		

\subsection{Criterio de estimaci\'on econ\'omica}

El an\'alisis de costes se ha realizado con enfoque de \textbf{coste de esfuerzo}.
La unidad base de c\'alculo es la hora de trabajo invertida por el equipo.

Se distinguen dos escenarios para evitar una lectura sesgada: el \textbf{Escenario A} (referencia acad\'emica) con tarifa de 15 EUR/h y el \textbf{Escenario B} (referencia profesional junior) con tarifa de 25 EUR/h. Esta doble referencia permite contextualizar el coste seg\'un perfiles de mercado diferentes.

La duraci\'on total considerada es de 660 horas.

\subsection{Coste de personal}

El coste principal del proyecto es el esfuerzo humano. Se calcula como:

\[
C_{personal} = H_{total} \times T_{hora}
\]

donde $H_{total}=660$ horas y $T_{hora}$ es la tarifa horaria del escenario.

\begin{table}[!htbp]
\centering
\begin{tabular}{|P{5.5cm}|P{2.8cm}|P{3.8cm}|}
\hline
\textbf{Escenario} & \textbf{Tarifa} & \textbf{Coste personal total} \\
\hline
Escenario A (acad\'emico) & 15 EUR/h & 9\,900 EUR \\
\hline
Escenario B (junior mercado) & 25 EUR/h & 16\,500 EUR \\
\hline
\end{tabular}
\caption{Coste de personal para 660 horas de proyecto}
\label{tab:coste-personal-total}
\end{table}

\subsection{Coste de personal por iteraci\'on}

Dado que los sprints regulares est\'an planificados en 120 horas y el sprint de cierre en 60 horas,
el coste por iteraci\'on regular es:

\[
C_{sprint} = 120 \times T_{hora}
\]


Para el sprint de cierre:

\[
C_{sprint\_cierre} = 60 \times T_{hora}
\]
\begin{table}[!htbp]
\centering
\begin{tabular}{|P{2cm}|P{2.2cm}|P{3.3cm}|P{3.3cm}|}
\hline
\textbf{Sprint} & \textbf{Horas} & \textbf{Coste A (15 EUR/h)} & \textbf{Coste B (25 EUR/h)} \\
\hline
S0 & 120 & 1\,800 EUR & 3\,000 EUR \\
\hline
S1 & 120 & 1\,800 EUR & 3\,000 EUR \\
\hline
S2 & 120 & 1\,800 EUR & 3\,000 EUR \\
\hline
S3 & 120 & 1\,800 EUR & 3\,000 EUR \\
\hline
S4 & 120 & 1\,800 EUR & 3\,000 EUR \\
\hline
S5 & 60 & 900 EUR & 1\,500 EUR \\
\hline
\textbf{Total} & \textbf{660} & \textbf{9\,900 EUR} & \textbf{16\,500 EUR} \\
\hline
\end{tabular}
\caption{Coste por sprint seg\'un escenario de tarifa}
\label{tab:coste-por-sprint}
\end{table}

\subsection{Coste asociado a aseguramiento de calidad}

La planificaci\'on incorpora tareas de testing en varias iteraciones
(unitarias, integraci\'on, E2E y mutaci\'on).
Tomando como referencia las horas de testing expl\'icitamente planificadas
en S1-S4, el esfuerzo de calidad asciende a:

\[
H_{QA} = 10 + 26 + 14 + 24 = 74\ \text{h}
\]

Esto representa un peso relativo de:

\[
\frac{74}{660} \times 100 = 11.21\%
\]

del esfuerzo total del proyecto, porcentaje coherente con un producto
que prioriza estabilidad antes de defensa.

\begin{table}[!htbp]
\centering
\begin{tabular}{|P{5.2cm}|P{2.2cm}|P{2.9cm}|P{2.9cm}|}
\hline
\textbf{Concepto} & \textbf{Horas} & \textbf{Coste A} & \textbf{Coste B} \\
\hline
Esfuerzo QA acumulado (S1-S4) & 74 & 1\,110 EUR & 1\,850 EUR \\
\hline
\end{tabular}
\caption{Coste estimado del aseguramiento de calidad}
\label{tab:coste-qa}
\end{table}

\subsection{Costes de infraestructura y operaci\'on}

Durante el desarrollo acad\'emico se han utilizado, de forma predominante,
planes gratuitos o cr\'editos de estudiante para servicios de integraci\'on,
despliegue y repositorio. Por tanto, el coste operativo \textbf{real observado}
ha sido pr\'oximo a 0 EUR en la etapa de desarrollo.

No obstante, para ofrecer una visi\'on profesional,
se incluye una estimaci\'on orientativa de operaci\'on m\'inima en un entorno no gratuito
durante 4 meses de proyecto:

\begin{table}[!htbp]
\centering
\begin{tabular}{|P{5.2cm}|P{3.2cm}|P{3.2cm}|}
\hline
\textbf{Concepto} & \textbf{Coste real (proyecto)} & \textbf{Coste orientativo (4 meses)} \\
\hline
Hosting backend & 0 EUR & 28 EUR \\
\hline
Base de datos gestionada & 0 EUR & 80 EUR \\
\hline
Servicios auxiliares de almacenamiento/transferencia & 0 EUR & 20 EUR \\
\hline
Dominio y gesti\'on b\'asica & 0 EUR & 4 EUR \\
\hline
\textbf{Total operaci\'on} & \textbf{0 EUR} & \textbf{132 EUR} \\
\hline
\end{tabular}
\caption{Comparativa entre coste real y coste operativo orientativo}
\label{tab:coste-operativo}
\end{table}

\subsection{Coste total del proyecto y escenarios de sensibilidad}

El coste total se modela como la suma del coste de personal y el coste operativo:
\[
C_{total} = C_{personal} + C_{operacion}
\]

Para ofrecer una visi\'on realista de la viabilidad del proyecto en distintos contextos, se plantea un an\'alisis de sensibilidad con tres escenarios de tarifas, manteniendo las 660 horas como constante de alcance. Se compara el coste asumiendo que no hay costes de operaci\'on (como ha ocurrido durante el desarrollo) frente a un escenario con operaci\'on orientativa de pago.

\begin{table}[!htbp]
\centering
\begin{tabular}{|P{3.5cm}|P{2.2cm}|P{3.5cm}|P{3.5cm}|}
\hline
\textbf{Escenario (Tarifa)} & \textbf{Coste personal} & \textbf{Total (Sin costes de operaci\'on)} & \textbf{Total (Con operaci\'on orientativa)} \\
\hline
Acad\'emico (15 EUR/h) & 9\,900 EUR & 9\,900 EUR & 10\,032 EUR \\
\hline
Junior mercado (25 EUR/h) & 16\,500 EUR & 16\,500 EUR & 16\,632 EUR \\
\hline
Perfil intermedio (35 EUR/h) & 23\,100 EUR & 23\,100 EUR & 23\,232 EUR \\
\hline
Perfil senior (45 EUR/h) & 29\,700 EUR & 29\,700 EUR & 29\,832 EUR \\
\hline
\end{tabular}
\caption{Coste total estimado del proyecto seg\'un escenario de tarifa}
\label{tab:coste-total-sensibilidad}
\end{table}

Este an\'alisis confirma que,
incluso con incrementos de tarifa,
la estructura modular y la automatizaci\'on de calidad mantienen
la viabilidad econ\'omica relativa del sistema frente al valor funcional entregado.

\subsection{Interpretaci\'on econ\'omica y sostenibilidad}

El an\'alisis muestra que el coste est\'a dominado por el esfuerzo de desarrollo,
mientras que el coste operativo es bajo en comparaci\'on.
Esto es coherente con la naturaleza de un producto software de tama\~no medio,
donde la inversi\'on principal est\'a en dise\~no, implementaci\'on, validaci\'on y documentaci\'on.

Adem\'as, la orientaci\'on modular del sistema reduce coste futuro de evoluci\'on,
ya que permite incorporar nuevas funcionalidades sin rehacer componentes nucleares.

Desde la perspectiva de sostenibilidad del producto,
la combinaci\'on de arquitectura desacoplada,
testing automatizado y despliegue continuo
reduce tambi\'en el coste de mantenimiento correctivo,
que suele representar una parte significativa del coste total de ciclo de vida.

\subsection{Conclusi\'on del an\'alisis de costes}

El coste global del proyecto resulta razonable para el alcance funcional conseguido
y para el nivel de calidad exigido en un TFG con vocaci\'on de producto real.
La estrategia de planificaci\'on por sprints, junto con el control de calidad continuo,
ha permitido mantener equilibrio entre plazo, alcance y robustez t\'ecnica.

	
